\documentclass[12pt]{article}
 \usepackage[margin=1in]{geometry} 
\usepackage{amsmath,amsthm,amssymb,amsfonts, enumitem, fancyhdr, color, comment, graphicx, environ}
\usepackage{booktabs}
\usepackage[colorinlistoftodos]{todonotes}
\usepackage{algorithm}
\usepackage{algpseudocode}
\usepackage{tikz}
\usepackage{subcaption}
\usepackage{float}
\usepackage{mathdots}
\usetikzlibrary{arrows.meta, automata, positioning}

\usepackage{listings}
\usepackage{xcolor}
\lstset{
  basicstyle=\ttfamily\small,
  backgroundcolor=\color{gray!10},
  frame=single,
  breaklines=true,
  keywordstyle=\color{blue},
  commentstyle=\color{gray},
  stringstyle=\color{red},
  showstringspaces=false,
  numbers=left,            
  numberstyle=\tiny\color{gray}, 
  numbersep=10pt 
}

%%%%%%%%%%%%%%%%%%%%%%%%%%%%%%%%%%%%%%%%%%%%%%%%%%%%%%%%%%%%%%%%%%%%%%%%
\pagestyle{fancy}
\setlength{\headheight}{65pt}
\newenvironment{problem}[2][Problem]{\begin{trivlist}
\item[\hskip \labelsep {\bfseries #1}\hskip \labelsep {\bfseries #2.}]}{\end{trivlist}}
\newenvironment{sol}
    {\emph{Solution:}
    }
    {
    \qed
    }
\specialcomment{com}{ \color{blue} \textbf{Comment:} }{\color{black}} %for instructor comments while grading
\NewEnviron{probscore}{\marginpar{ \color{blue} \tiny Problem Score: \BODY \color{black} }}

\newcounter{subproblem}
% \renewcommand{\thesubproblem}{\alph{subproblem}} % letters 
\renewcommand{\thesubproblem}{\arabic{subproblem}} % numbers
\newenvironment{subprob}[1][]{
  \refstepcounter{subproblem}
  \begin{trivlist}
  \item[\hskip \labelsep {\bfseries (\thesubproblem)}]
}{
  \end{trivlist}
}
\newenvironment{subsol}
    {\emph{Solution:}
    }
    {
    \qed
    }

\newtheorem{theorem}{Theorem}
\newtheorem{lemma}{Lemma}
\newtheorem{proposition}{Proposition}
\newtheorem{corollary}{Corollary}

% numbered within sections (optional)
% \newtheorem{theorem}{Theorem}[section]
% \newtheorem{lemma}[theorem]{Lemma} % shares theorem counter

%%%%%%%%%%%%%%%%%%%%%%%%%%%%%%%%%%%%%%%%%%%%%%%%%%%%%%%%%%%%%%%%%%%%%%%%%%%%%%%%%
\setlength {\marginparwidth }{2cm}
\newcommand{\nextline}{\\[0.2em]}


\newtheorem*{remark}{Remark}

\newcommand{\R}{\mathbb{R}}
\newcommand{\N}{\mathbb{N}}
\newcommand{\Q}{\mathbb{Q}}
\newcommand{\Z}{\mathbb{Z}}
\newcommand{\st}{\text{ s.t }}
\newcommand{\tand}{\text{ and }}
\newcommand{\tor}{\text{ or }}
\newcommand{\bigO}[1]{\mathcal{O}\left(#1\right)}
\newcommand{\fP}{\mathbb{P}}
\newcommand{\E}{\mathbb{E}}

%%%% In most cases you won't need to edit anything above this line %%%%

\begin{document}

\hfill CMSC 27100 Practice Problems

\hfill Mideterm Review

\section*{Logic}
\begin{enumerate}
  \item Determine if the statement is true or false. If false, give a counterexample. 
  \[\forall a,b \in \R, ~a|b \rightarrow a^2|b^2\]

  \item Prove or Disprove: If $p\rightarrow q$ and $q \rightarrow r$, then $(p \vee q) \rightarrow r$
\end{enumerate}


\section*{Cartesian Products}
\begin{enumerate}
  \item If $|A| = n \tand |B| = m$, prove $|A \times B| = nm$
  \item How many subsets of $|A \times B|$ exist? Assume $|A| = n, |B| = m, \tand |A \cap B| = 0$.
\end{enumerate}

\section*{Set Theory}
\begin{enumerate}
  \item Prove that $(A \setminus B) \cap B = \emptyset$
  \item Assume $|A| = 10, |B| = 30, \tand |A \cap B| = 10$. What can we infer about $A \tand B$?
\end{enumerate}

\section*{Functions}
\begin{enumerate}
  \item If $G: A \rightarrow B$ and $F: B \rightarrow C$, prove $F \circ G: A \rightarrow C$ is injective.
  \item If $F: A \rightarrow B$ is bijective, prove that $F^{-1}$ exists and is bijective.
\end{enumerate}

\section*{Number Theory}
\begin{enumerate}
  \item Prove that $101$ is prime.
  \item Prove that if $\gcd(a, b) = 1$, then $a|bc \rightarrow a|c$
\end{enumerate}

\section*{Division Theorem}
\begin{enumerate}
  \item Write $-57$ in the form: $8q + r$ with $0 \leq r < 8$
  \item Prove uniqueness of the division theorem. (Existance was previously proven in class)
\end{enumerate}

\section*{Euclid's Algorithm}
\begin{enumerate}
  \item Find integers $x, y$ s.t. $414x + 662y = \gcd(414,662)$
  \item Find multiplicative inverse of $37 \pmod{101}$
\end{enumerate}

\section*{Modular Arthemtic and Fermat's Little Theorem}
\begin{enumerate}
  \item Determine if $18$ has an inverse modulo $30$
  \item Compute $7^{100} \pmod{13}$
  \item Prove that if $p$ is prime and $p \nmid a$, then $a^{p-2} = a^{-1} \pmod{p}$
\end{enumerate}

\section*{Equivalence Classes and CRT}
\begin{enumerate}
  \item How many equivalence classes are there under mod $n$ on $\Z$?
  \item Find integer $x$ s.t. $x \equiv 1 \pmod{4} \tand x \equiv 3 \pmod{6}$. Solve or explain why it is impossible.

\end{enumerate}

\section*{Induction}
\begin{enumerate}
  \item Prove that any integer $n \geq 8$ can be made of a combonation of units of $5$ and $3$.
  \item We define the fibonacci sequence s.t. $n_1 = 0, n_2 = 1 \tand \forall n > 2, a_n = a_{n-1} + a_{n-2}$. Prove that $a_n < 2^n ~\forall n \geq 1$.
\end{enumerate}

\section*{Wilson's Theorem}
\begin{enumerate}
  \item Prove both direction of Wilson's Theorem. In other words, prove that $p$ is prime $\iff (p-1)! \equiv -1 \pmod{p}$
\end{enumerate}



\end{document}